% Part: set-theory
% Chapter: ordinals
% Section: introduction

\documentclass[../../../include/open-logic-section]{subfiles}

\begin{document}

\olfileid[id]{sth}{ordinals}{intro}

\olsection{Pendahuluan}

Dalam \olref[z][]{chap}, kita mempostulatkan adanya tahap tak hingga dalam
hierarki melalui \stagesinf{} (lihat pula Aksioma Ketakhinggaan kita).
Namun, dengan \stagessucc{}, kita tidak dapat berhenti pada tahap tak hingga;
kita harus terus melanjutkan. Jadi, pada tahap sesudah tahap
tak hingga pertama, kita membentuk semua kumpulan himpunan yang mungkin dari
himpunan-himpunan yang tersedia pada tahap tak hingga pertama; lalu
mengulanginya lagi, lagi, lagi, \dots

Secara implisit, yang terjadi di sini ialah kita mulai menggunakan gagasan
``intuitif'' tentang bilangan, yang menurutnya dapat ada bilangan
\emph{sesudah} semua bilangan asli. Secara khusus, gagasan yang terlibat ialah
gagasan \emph{ordinal transfinit}. Tujuan bab ini adalah membuat gagasan
tersebut lebih teliti. Kita akan menjelajahi gagasan umum ordinal, lalu secara
eksplisit mendefinisikan himpunan-himpunan tertentu sebagai ordinal kita.

\end{document}
